\documentclass[11pt,a4paper]{article}
\usepackage[margin=1in]{geometry}
\usepackage{amsmath,amssymb,amsthm}
\usepackage{graphicx}
\usepackage{microtype}
\usepackage[hidelinks]{hyperref}

\newtheorem{theorem}{Theorem}
\newtheorem{proposition}[theorem]{Proposition}
\newtheorem{remark}[theorem]{Remark}
\newcommand{\N}{\mathbb N}
\newcommand{\one}{\mathbf 1}
\newcommand{\supp}{\operatorname{supp}}

\title{An un-null net with no un-null sequence of its terms}
\author{A counterexample to Question 2.14 of arXiv:1605.03538}
\date{17 August 2026}

\begin{document}
\maketitle

\noindent\textbf{Status.}
Candidate counterexample, likely valid, pending human review.  Question~2.14
has a negative answer.  In fact, there is a positive normalized un-null net
from which \emph{no} sequence of terms is un-null, even if the indices need
not be increasing.

\section{Source question}

For a Banach lattice $X$, a net $(x_\alpha)$ is \emph{un-null} if
\[
  \bigl\|\,|x_\alpha|\wedge u\,\bigr\|\longrightarrow0
  \qquad(u\in X_+).
\]
Corollary~2.13 of the source proves that, when $X$ has a quasi-interior
point, every un-null net contains an increasing sequence of indices whose
terms are still un-null.  Question~2.14 on source PDF page~5 asks whether
this remains true without a quasi-interior point~\cite{DOT}.

\begin{center}
\includegraphics[width=\textwidth]{figures/open_problem_crop.png}
\end{center}

\section{Proof intuition}

The failure comes from a clash between two countabilities.  Every element of
$\ell_\infty^c(\omega_1)$ sees only countably many coordinates, so a net can
eventually escape the support of each fixed test vector.  But after anyone
chooses countably many terms of the net, all coordinates used by those terms
again form one countable set.  Its indicator is a single legal test vector
which detects every chosen term with norm one.

\section{The counterexample}

Let $\Gamma=\omega_1$, the first uncountable ordinal, and let
\[
 X=\ell_\infty^c(\Gamma)
   :=\{x\in\ell_\infty(\Gamma):\supp x\text{ is countable}\}
\]
with the coordinatewise order and the supremum norm.  Let
$\mathcal A=[\Gamma]^{\leq\aleph_0}$ be the directed set of countable subsets
of $\Gamma$, ordered by inclusion.  For $C\in\mathcal A$, put
\[
  \gamma_C=\min(\Gamma\setminus C),
  \qquad x_C=e_{\gamma_C}.
\]
The minimum exists because $C$ is countable whereas $\Gamma$ is not.

\begin{theorem}\label{thm:main}
The space $X$ is a Banach lattice with no quasi-interior point, and the net
$(x_C)_{C\in\mathcal A}$ is positive, normalized, and un-null.  Nevertheless,
for every sequence $(C_n)$ in $\mathcal A$ there is a single $u\in X_+$ such
that
\[
  \bigl\|\,|x_{C_n}|\wedge u\,\bigr\|=1
  \qquad(n\in\N).
\]
Consequently, no sequence of terms of this net is un-null; in particular,
Question~2.14 has a negative answer.
\end{theorem}

\begin{proof}
The space $X$ is a vector sublattice and an ideal of $\ell_\infty(\Gamma)$.
It is norm closed: if $z_k\in X$ and $z_k\to z$ uniformly, then
\[
  \supp z\subseteq\bigcup_{k=1}^\infty\supp z_k,
\]
which is countable.  Thus $X$ is a Banach lattice.

No $e\in X_+$ is quasi-interior.  Indeed, $\supp e$ is countable, so one may
choose $\eta\in\Gamma\setminus\supp e$.  Every vector in the ideal generated
by $e$, and hence every vector in its norm closure, vanishes at $\eta$;
the vector $e_\eta$ does not.  Thus the closed ideal generated by $e$ is
proper.

Now fix $u\in X_+$ and put $D=\supp u$.  This set is countable and hence is
an index in $\mathcal A$.  Whenever $C\geq D$, the definition gives
$\gamma_C\notin C\supseteq D$.  Therefore
\[
  |x_C|\wedge u=e_{\gamma_C}\wedge u=0.
\]
This holds eventually for every $u\in X_+$, so $x_C\to0$ in the un-topology.
Also $x_C\geq0$ and $\|x_C\|_\infty=1$ for every $C$.

Finally, let $(C_n)$ be any sequence of indices and set
\[
  E=\{\gamma_{C_n}:n\in\N\},\qquad u=\one_E.
\]
The set $E$ is countable, so $u\in X_+$.  Since
$u(\gamma_{C_n})=1$, we have
\[
  |x_{C_n}|\wedge u=e_{\gamma_{C_n}}
  \quad\text{and hence}\quad
  \bigl\|\,|x_{C_n}|\wedge u\,\bigr\|_\infty=1
\]
for every $n$.  Thus $(x_{C_n})$ is not un-null.  The argument does not use
that $(C_n)$ is increasing, proving the stronger assertion.
\end{proof}

\section{Strength, checks, and literature boundary}

The initial search through $c_0$-sums failed because their uniformly small
tails permit a diagonal sequence extraction.  Replacing $c_0$ support by
arbitrary countable support is the decisive upgrade: fixed un-tests remain
countable, while the indicator of all coordinates used by a proposed
sequence remains in the lattice.  The example is therefore stronger than
required: the net consists of positive norm-one atoms, and no sequence of
its values is un-null.  The finite incidence identities underlying the two
support arguments are independently exercised in
\texttt{code/verify\_incidence.py}; the infinite conclusion itself is the
set-theoretic proof above.

A bounded primary-source search used the exact question, un-null nets,
embedded sequences, the un-topology, and countably supported bounded
functions.  The 2023 paper of Deng and de Jeu proves that every un-convergent
net has an embedded sequence that is uo-convergent, and obtains a sequence
that is also un-convergent under order continuity~\cite{DJ}; it does not
state the unrestricted conclusion refuted here.  No primary arXiv source
stating this counterexample was found.  Correctness confidence is high;
novelty confidence is moderate pending a broader expert search.

The source contains a second open question: whether weak nullity of every
norm-bounded positive un-null net forces $X^*$ to be order continuous.  That
question was subsequently answered affirmatively by Elbour's
Theorem~2.16~\cite{Elbour}.  Indeed, the positive-net hypothesis applied to
$|x_\alpha|$ makes every functional vanish on every norm-bounded un-null net,
and the cited theorem with scalar target yields order continuity of $X^*$.
This literature fact is separate from Theorem~\ref{thm:main}.

\begin{thebibliography}{3}
\bibitem{DOT}
Y.~Deng, M.~O'Brien, and V.~G.~Troitsky,
\emph{Unbounded norm convergence in Banach lattices}, arXiv:1605.03538.

\bibitem{DJ}
Y.~Deng and M.~de Jeu,
\emph{Embedded unbounded order convergent sequences in topologically
convergent nets in vector lattices}, arXiv:2309.03027.

\bibitem{Elbour}
A.~Elbour,
\emph{Some properties on the unbounded absolute weak convergence in Banach
lattices}, arXiv:2004.10691.
\end{thebibliography}

\end{document}
